\scanpage{99}
Comme $\mathcal A_i\simeq E_i\otimes L_i^\vee$, $\mathcal A\simeq E\otimes L^\vee$, on a
\[
E\otimes L^\vee\longleftrightarrow(E_1\otimes L_1^\vee)\otimes(E_2\otimes L_2^\vee)
\]
et comme $L^\vee\simeq L_1^\vee\otimes L_2^\vee$, on en déduit également
\[
E\xrightarrow{i}E_1\otimes E_2
\]
qui est une sorte de transposé de $E_1\otimes E_2\longrightarrow E$.
\margin{$x_1\otimes x_2\mapsto x_1*x_2$;}
\deleted{Pour expliciter $i$, il faut donner les deux composés
$\inserted{\alpha=}i\circ *$, $\inserted{\beta=}i\circ\mathrm{inj}$,
en sorte que $\alpha\circ\mathrm{inc}=2\beta$.}
\deleted{\diagram{\begin{tikzcd}[column sep=large]
E_1\otimes E_2 \arrow[dr,"*"] \arrow[rr,bend left=25,"?"] & & E_1\otimes E_2\\
\begin{array}{c}L_1\otimes E_2\oplus E_1\otimes L_2\\E_1*_0E_2\end{array} \arrow[u,hook,"\mathrm{inc}"] \arrow[drr,no head,"\deleted{\unsure{\mathrm{proj}}}"'] & E \arrow[ur,"i"] & \\
& L_1\otimes L_2 \arrow[u,hook,"\mathrm{inj}"] \arrow[uur,bend right=25,"?\ \text{i.e. can.}"'] &
\end{tikzcd}}}
\deleted{\deleted{Je pense} on voudra $L_1\otimes L_2\longrightarrow E_1\otimes E_2$ inclusion canonique, donc un \unsure{hom}}
\deleted{\diagram{\begin{tikzcd}
E_1\otimes_0 E_2\;\deleted{\unsure{L_2}} \arrow[r,no head] & \unsure{E_1\otimes L_2}\subset E_1\otimes E_2\\
E_1\otimes E_2 \arrow[u,no head,"\unsure{\cap}"] \arrow[ur,dashed,"?"'] &
\end{tikzcd}}}
\deleted{(pour $\boxed{\uuline{\unsure{\chi}=2.}}$),}
\deleted{\deleted{Je pense} Considérons les inv. canoniques $\sigma_1,\sigma_2$ de $E_1$, de $E_2$, d'où $\sigma=\sigma_1\otimes\sigma_2$ de $\unsure{E}$, noté $\xi\mapsto\bar\xi$. Je pense à l'hom
\[
\xi\longmapsto\xi+\bar\xi
\]
donc
\[
\begin{aligned}
x_1\otimes x_2&\longmapsto x_1\otimes x_2+\bar x_1\otimes\bar x_2\\
&=x_1\otimes x_2+(\unsure{\overline T_1}\inserted{\psi_1(x_1)}-x_1)\otimes(\unsure{\overline T_2}\inserted{\psi_2(x_2)}-x_2).
\end{aligned}
\]
Si $x_1=u_1\in\Gamma L_1$, alors on trouve $\bar x_1=-u_1$,
\[
u_1\otimes x_2\longmapsto u_1\otimes x_2+(-u_1)(\unsure{\overline T_2}\psi_2(x_2)-x_2)
\]
}
\mspage{98}
